\documentclass{article}
\usepackage[utf8]{inputenc}
\usepackage{plcourse}
\usepackage{ttquot}
\usepackage{proof}
\usepackage{stmaryrd}
\usepackage{angle}
\usepackage{amssymb}
\usepackage{hyperref}
\usepackage{graphicx}
\usepackage{tikz}
\usepackage{tikz-qtree}
\usepackage{xspace}


\homework{5}

\begin{document}
\hwsubheader


% %%%%%%%%%%%%%%%%%%%%%%%%%%
%
%    Cover page
%
%
\vspace{5cm}
{\LARGE
\begin{tabular}{rp{0.6\linewidth}}
  Name:&Pranav Sitaraman\\
  HUID:&31712165\\
  Collaborators:&{\normalsize
    Stephen Gwon, Gavin Ye.
    }
\end{tabular}
}

\vfill
\textit{Make sure that the remaining pages of this assignment do not contain any identifying information.}
\vfill

\newcommand{\bindop}{\mathbin{\mbox{\ttfamily >>=}}}
\newcommand{\sappend}{\mathbin{++}}
\newcommand{\emptyout}{\epsilon}


% %%%%%%%%%%%%%%%%%%%%%%%%%%
%
%    Question 1
%
%


\newpage

\begin{question}{Monad Laws}{(25 points)}

The useful way to read the \texttt{Writer} bind is that it runs the first
computation, runs the second computation using the first value, and then
keeps the second value while appending the two pieces of output in order.
Thus the three monad laws should come exactly from the three monoid laws for
\texttt{mappend}.

\medskip
\noindent\textbf{Left unit.}
Suppose \(f~x = \code{Writer}(y,w_f)\). Then
\[
\begin{aligned}
  \code{return}\ x \bindop f
    &= \code{Writer}(x,\code{mempty}) \bindop f \\
    &= \code{Writer}(y,\code{mappend}\ \code{mempty}\ w_f) \\
    &= \code{Writer}(y,w_f) \\
    &= f~x.
\end{aligned}
\]
The only real step is the third one, which is the left-unit law for the
monoid \(w\): \(\code{mappend}\ \code{mempty}\ w_f=w_f\).

\medskip
\noindent\textbf{Right unit.}
Let \(xM=\code{Writer}(x,w)\). Then
\[
\begin{aligned}
  xM \bindop \code{return}
    &= \code{Writer}(x,w) \bindop \code{return} \\
    &= \code{Writer}(x,\code{mappend}\ w\ \code{mempty}) \\
    &= \code{Writer}(x,w) \\
    &= xM.
\end{aligned}
\]
Here we used the right-unit law
\(\code{mappend}\ w\ \code{mempty}=w\).

\medskip
\noindent\textbf{Associativity.}
Let
\[
  xM=\code{Writer}(x,w_1), \qquad
  f~x=\code{Writer}(y,w_2), \qquad
  g~y=\code{Writer}(z,w_3).
\]
Expanding the left side first,
\[
\begin{aligned}
  (xM \bindop f)\bindop g
    &= (\code{Writer}(x,w_1)\bindop f)\bindop g \\
    &= \code{Writer}(y,\code{mappend}\ w_1\ w_2)\bindop g \\
    &= \code{Writer}
        (z,\code{mappend}\ (\code{mappend}\ w_1\ w_2)\ w_3).
\end{aligned}
\]
On the other hand,
\[
\begin{aligned}
  xM \bindop (\lambda x.\ f~x\bindop g)
    &= \code{Writer}(x,w_1)\bindop(\lambda x.\ f~x\bindop g) \\
    &= \code{Writer}(x,w_1)\bindop
        (\lambda x.\ \code{Writer}
          (z,\code{mappend}\ w_2\ w_3)) \\
    &= \code{Writer}
        (z,\code{mappend}\ w_1\ (\code{mappend}\ w_2\ w_3)).
\end{aligned}
\]
These two final outputs are equal by associativity of the monoid operation:
\[
  \code{mappend}\ (\code{mappend}\ w_1\ w_2)\ w_3
  =
  \code{mappend}\ w_1\ (\code{mappend}\ w_2\ w_3).
\]
The value part is \(z\) on both sides, so the two \texttt{Writer} computations
are equal. Therefore \(\code{Writer}\ w\) satisfies all three monad laws
whenever \(w\) is a monoid.

\end{question}

% %%%%%%%%%%%%%%%%%%%%%%%%%%
%
%    Question 2
%
%

\newpage

\begin{question}{Lambda-Print}{(25 points)}

I will write \(\emptyout\) for the empty output string, and
\(\sappend\) for semantic string append. The rules below are deliberately
left-to-right: in an application we evaluate the function, then the argument,
then the body; in a \(\code{concat}\) we evaluate the left expression before
the right one. I also use \(e[v/x]\) for the usual capture-avoiding
substitution of value \(v\) for \(x\) in \(e\).

\[
\infrule
  {}
  {\lam{x}{e} \Downarrow \langle \lam{x}{e},\emptyout\rangle}
  {\textsc{Lam}}
\qquad
\infrule
  {}
  {s \Downarrow \langle s,\emptyout\rangle}
  {\textsc{String}}
\qquad
\infrule
  {}
  {\code{unit} \Downarrow \langle \code{unit},\emptyout\rangle}
  {\textsc{Unit}}
\]

\[
\infrule
  {e_1 \Downarrow \langle \lam{x}{e},o_1\rangle
   \qquad
   e_2 \Downarrow \langle v_2,o_2\rangle
   \qquad
   e[v_2/x] \Downarrow \langle v,o_3\rangle}
  {e_1~e_2 \Downarrow \langle v,o_1\sappend o_2\sappend o_3\rangle}
  {\textsc{App}}
\]

\[
\infrule
  {e_1 \Downarrow \langle s_1,o_1\rangle
   \qquad
   e_2 \Downarrow \langle s_2,o_2\rangle}
  {\code{concat}\ e_1\ e_2
   \Downarrow
   \langle s_1\sappend s_2,o_1\sappend o_2\rangle}
  {\textsc{Concat}}
\qquad
\infrule
  {e \Downarrow \langle s,o\rangle}
  {\code{output}\ e
   \Downarrow
   \langle \code{unit},o\sappend s\rangle}
  {\textsc{Output}}
\]

There is no separate rule for a free variable. In this substitution-style
semantics, closed programs only evaluate variables after an application has
substituted an argument value for the bound variable. Similarly,
\(\code{concat}\) and \(\code{output}\) only have rules when their evaluated
arguments have string values; otherwise there is simply no large-step
derivation.

\medskip
\noindent
As a quick check, for string literals \(s_1,s_2\),
\[
\begin{aligned}
  \code{output}\ (\code{concat}\ s_1\ s_2)
  &\Downarrow
    \langle \code{unit},\emptyout\sappend(s_1\sappend s_2)\rangle \\
  &= \langle \code{unit},s_1\sappend s_2\rangle.
\end{aligned}
\]
For the expression \((\lambda x.\ \code{output}\ s_2)\;(\code{output}\ s_1)\),
the function itself produces no output, the argument produces \(s_1\), and
the body produces \(s_2\). The application rule therefore gives
\[
  (\lambda x.\ \code{output}\ s_2)\;(\code{output}\ s_1)
  \Downarrow
  \langle \code{unit},\emptyout\sappend s_1\sappend s_2\rangle
  =
  \langle \code{unit},s_1\sappend s_2\rangle.
\]
So the two programs produce the same value and output, as required. The
example from the problem statement is the same calculation with
\(s_1=\text{``ab''}\), \(s_2=\text{``cd''}\), followed by the final output
\(\text{``ef''}\), giving output \(\text{``abcdef''}\).

\end{question}


% %%%%%%%%%%%%%%%%%%%%%%%%%%
%
%    Question 3
%
%


\newpage

\begin{question}{Monadic Interpreter}{(Optional)}

I completed the interpreter in \texttt{eval.hs}. The implementation follows
the rules above, but uses environment semantics instead of explicit
substitution: evaluating a lambda produces a closure with the current
environment, and applying a closure evaluates the body in the closure
environment extended with the argument value.

The monad representation is
\[
  \code{MaybeWriter}\ w\ a
  =
  \code{MaybeWriter}\;(\code{Maybe}\ (a,w)),
\]
so a successful computation stores both its value and accumulated output,
while failure is represented by \(\code{Nothing}\). The bind operation
sequences computations and appends the outputs in the same order as the
large-step rules:
\[
  \code{Just}\ (x,w_1) \bindop f
  \quad\text{runs } f~x=\code{Just}\ (y,w_2)
  \quad\text{and returns}\quad
  \code{Just}\ (y,w_1\sappend w_2).
\]

The evaluator itself stays in monadic style: variable lookup, closure
checking, and string checking are small helper computations that either
\(\code{return}\) the coerced value or produce \(\code{mwFailure}\). For
example, the output case is exactly the operational rule:
\begin{verbatim}
evalM (Output e) env = do
  v <- evalM e env
  s <- asString v
  mwOutput s
  return UnitVal
\end{verbatim}
Finally, \(\code{eval}\) simply runs the monadic computation and exposes the
pure result of type \(\code{Maybe}\ (\code{Val},\code{String})\), which is the
format requested by the assignment.

\end{question}


\end{document}
